\section{From polydiscs to convex approximation}\label{sec:globalization}

We now have approximation on polydiscs in every dimension. Convex
sets need not fit a polydisc inside the original map domain, so this
alone does not prove CAP. The required enlargement has the form
\[
 Q_-=Q\cap\{\operatorname{Im}z_1\leq0\}\subset Q,
\]
where $Q$ is compact and convex. Given a map near $Q_-$, we will
first approximate it on the joining real face by an entire map,
with control of tangential derivatives. That map supplies values on
the upper side. Its accuracy need not persist on any fixed complex
collar of the face, so the second step is a gluing estimate measured
on the face itself. Repeating this enlargement across the faces of
a polytope will prove Theorem~\ref{ana:cap}.

The induction must retain arbitrary extra polydisc factors: they
include both genuine source parameters and the auxiliary variables
of the chart sprays used to glue maps.

For \(r\geq0\), denote by \({\cal A}_r\) the following statement:
for every compact convex \(C\subset\C^r\), every extra closed
polydisc \(P\), and every map holomorphic near \(C\times P\),
there are entire maps approximating it uniformly on \(C\times P\).
All the approximations in \({\cal A}_r\) may also be required in
any fixed finite number of derivatives on a slightly smaller
compact product.  To justify this last formulation, first
thicken \(C,P\) within the given domain, approximate uniformly
there, and use chart-spray coordinates and Cauchy estimates on
the smaller product.  This uses no new approximation assertion.

\subsection{Holomorphic approximation with a real parameter}

The slices of a convex real face vary with its real coordinate.
We will join approximations on neighboring slices by a smooth
real-parameter path. The next lemma makes such a path holomorphic
near the real interval; the following lemma then makes it entire.

\begin{lemma}[Complexifying a real parameter]\label{glue:interval}
Let $I\subset\R$ be a compact interval and let $C\subset\C^k$ be a
compact convex set. Suppose $H_s(q)$ is a smooth family, for $s$ in a
neighborhood of $I$, of holomorphic maps on a common neighborhood of
$C$, with values in a complex surface. Smoothness is understood on a
common neighborhood with all finite derivatives under consideration.
For every finite $j$ and every $\epsilon>0$, after an arbitrarily small
loss of the specified compact margins, $H$ can be approximated in
$C^j(I\times C)$ by a map holomorphic on a neighborhood of $I\times C$
in $\C\times\C^k$. The size of this complex neighborhood may depend on
the approximant. Additional closed polydisc factors can be included in
$C$, with fixed larger neighborhoods.
\end{lemma}

\begin{proof}
We give the gluing detail to avoid assuming a holomorphic extension of
the original smooth family. By Lemma~\ref{loc:spray}, a finite cover of
$I$ by short intervals admits fixed graph charts
$E_i(q,v)$ over the corresponding maps. First choose a smooth family of
small chart sprays $H_s^\sharp(q,v)$ centered at $H_s(q)$, holomorphic in
$(q,v)$, on a common small fiber ball. This can be done explicitly along
the interval. In one fixed chart, translate its fiber coordinates to
the center representing $H_s$. When changing charts at a division point
$s_0$, express the preceding spray in the new chart and take its fiber
displacement at $s_0$. On a short overlap interpolate this displacement
to the fixed displacement at $s_0$, leaving the moving center unchanged.
The two displacements are as close in the first fiber derivative as
desired when the overlap is short. The derivatives therefore remain
invertible. Finitely many such operations and a smaller common fiber
ball give the asserted smooth spray family.

In each division interval write this family in one fixed graph chart.
Approximate its coordinate functions by polynomials in $s$ with
coefficients holomorphic in $(q,v)$, retaining any prescribed finite
number of derivatives, and interpolate its endpoint jets through order
$M$. This is ordinary approximation in a Banach space of holomorphic
functions on a smaller fixed neighborhood: approximate a sufficiently
high real derivative uniformly, integrate, and make a finite Hermite
correction at the endpoints. The correction tends to zero with the
initial approximation error. The target sprays obtained on adjacent
intervals consequently have identical $s$-jets through order $M$ at the
division point.

Now, after the finitely many polynomials have been fixed, choose a small
complex thickness $\delta>0$ around the real intervals. Extend each
interval slightly past its endpoints, so adjacent rectangles overlap
within distance $O(\delta)$ of their common division point. Their
spray transition differs from the identity by $O(\delta^{M+1})$ there,
uniformly on a fixed smaller fiber ball and parameter compactum. This
is Taylor's formula applied to the holomorphic coordinate transitions;
the constants are fixed before $\delta$ is chosen.

Additive splitting on these finitely many thin planar rectangles costs
at most $C\delta^{-1}$. Indeed, use cutoffs across the vertical overlap
bands, whose first derivatives are $O(\delta^{-1})$, and the
Cauchy--Green operator on a bounded planar set. For simultaneous
splitting, each cut contributes a pair of corrections on the entire
portion to its left and right; that contribution has zero jump at every
other cut. The number of cuts is fixed. The source variables $(q,v)$
are passive for this operator, so no parameter radius is lost.

For clarity, write the consecutive transitions as
$\gamma_i=\operatorname{id}+c_i$, with
$\|c_i\|+\|D_vc_i\|\le A\delta^{M+1}$ on a fixed smaller fiber ball.
Let $R$ be the simultaneous right inverse, of norm at most $K/\delta$,
for the jump operator $(a_i)\mapsto(a_{i+1}-a_i)$.
The central maps are identified by the fixed-point equation
\[
                 (a_i)=R\bigl((c_i(\,\cdot\,,a_i))_i\bigr).
\]
It preserves the ball of radius $2KA\delta^M$ and is contractive there
when $KA\delta^M<1/2$. The corrections are therefore $O(\delta^M)$.
Substituting them into the corresponding sprays gives central maps
agreeing on overlaps. Cauchy estimates on rectangles with margins
comparable to $\delta$ bound their correction derivatives through order
$j$ by $O(\delta^{M-j})$. Derivatives in the other variables use their
fixed margins. Take $M>j+2$ and then $\delta$ small. Together with the
initial polynomial accuracy this gives the prescribed $C^j$ estimate.
\end{proof}
\begin{lemma}\label{ana:parameter-disc}
Assume \({\cal A}_r\).  Let \(I\subset\R\) be a compact interval,
\(C\subset\C^r\) compact and convex, and \(P\) a closed polydisc.
A map holomorphic near \(I\times C\times P\) can be approximated
there, with any prescribed finite derivative accuracy and smaller
positive margins, by entire maps of all the displayed complex
variables.
\end{lemma}

\begin{proof}
First shrink any unused margins so that the given map \(H\) is
defined on \(V\times U\), where \(V\subset\C\) is a simply
connected bounded neighborhood of \(I\), and \(U\) contains a
slightly thickened compact product \(C\times P\).
A thin open rectangle is a possible \(V\).
Let \(\phi:V\to\Delta\) be a Riemann map.  Choose \(0<\rho<1\)
so that \(\phi(I)\Subset\Delta_\rho\).
The map
\[
 (\zeta,y,u)\longmapsto H(\phi^{-1}(\zeta),y,u)
\]
is defined near
\(\overline\Delta_\rho\times C\times P\), with slightly smaller
positive margins.  Apply \({\cal A}_r\), using the \(\zeta\)
disc as an extra passive factor.  This gives an entire map
\(E(\zeta,y,u)\) with arbitrarily good derivative approximation
on a smaller product still containing \(\phi(I)\times C\times P\).
Consequently \(E(\phi(s),y,u)\) approximates \(H(s,y,u)\) with
the required derivatives along \(I\times C\times P\).

Choose a closed convex neighborhood \(I'\) of \(I\) contained in
\(V\).  Scalar polynomial approximation, with Cauchy estimates
after a further harmless shrink, approximates \(\phi\) and any
specified finite number of its derivatives on \(I'\) by a
polynomial \(p\).  At this stage \(E\) is fixed.  On the compact
sets in question all its derivatives up to the required order
are bounded.  The chain rule therefore shows that
\(E(p(s),y,u)\) is as close to \(E(\phi(s),y,u)\) as desired.
The map \(E(p(s),y,u)\) is entire in every variable.
\end{proof}

\subsection{Approximation on a convex real face}

\begin{lemma}\label{ana:face}
Assume \({\cal A}_{n-1}\), where \(n\geq1\), and
Proposition~\ref{ana:polydisc}.  Let
\(F\subset\R\times\C^{n-1}\) be compact and convex, and let \(P\)
be a closed polydisc.  Every holomorphic map near \(F\times P\)
can be approximated on \(F\times P\) by entire maps on
\(\C^n\times\C^{\dim P}\), with any fixed finite derivative
accuracy along the real face and smaller positive margins.
\end{lemma}

\begin{proof}
Write the variables as \((t,y,u)\), with \(t\) real on the face.
Thicken \(F\) slightly within \(\R\times\C^{n-1}\), and thicken
\(P\), so that all eventual estimates take place strictly inside
the initial domain.  Denote the thickened face again by \(F\).
Its projection is a compact interval \(I\), and write
\(F_t=\{y:(t,y)\in F\}\).

We first specify the finite cover used below.  For each \(t_0\in I\),
choose a compact convex neighborhood \(C_{t_0}\) of \(F_{t_0}\)
in the \(y\) variables and an open complex disc \(D_{t_0}\) about \(t_0\)
such that the given map \(f\) is defined near
\[
 \overline D_{t_0}\times C_{t_0}\times P,
\]
including some fixed smaller margins.  These choices are possible
by compactness of \(F_{t_0}\times P\).  After reducing the real
interval about \(t_0\), it has
\[
 F_t\Subset\operatorname{int}C_{t_0}
\]
uniformly on that interval.  Indeed, otherwise a convergent
sequence \((t_\nu,y_\nu)\in F\), with \(t_\nu\to t_0\), would
contradict \(F_{t_0}\subset\operatorname{int}C_{t_0}\).
Choose a finite subdivision of \(I\), with small overlaps at
its division points, subordinate to these intervals.
Denote the corresponding products by \(D_i\times C_i\times P\).
All sets can be chosen a little larger than the sets used in
the subsequent estimates.

At a division point \(b\), both adjacent \(C_i\)'s contain
\(F_b\) in their interiors.  Choose compact convex sets \(C^-\)
and \(C\) with nonempty interiors such that
\[
 F_b\subset\operatorname{int}C^-,
 \qquad C^-\Subset\operatorname{int}C,
 \qquad C\Subset\operatorname{int}(C_i\cap C_{i+1}).
\]
By the preceding compactness argument, after reducing a complex
disc \(D_b\) about \(b\), the sets \(F_t\), for real \(t\) near
\(b\), lie strictly inside \(C^-\), and
\[
 \overline D_b\times C\times P
       \Subset (D_i\times\operatorname{int}C_i\times
                    \operatorname{int}P^+)
          \cap(D_{i+1}\times\operatorname{int}C_{i+1}\times
                    \operatorname{int}P^+)
\]
for an available larger passive polydisc \(P^+\).
All transition intervals and their cutoff functions are fixed now.
Also fix a large \(y\)-polydisc \(B\) containing all the \(C_i\)
and all face projections, with a positive margin.

By \({\cal A}_{n-1}\), applied with the \(t\) disc and \(u\)
polydisc passive, choose entire maps \(g_i(t,y,u)\) arbitrarily
close to \(f\), with any fixed finite derivative accuracy, on
\(\overline D_i\times C_i\times P\).
More explicitly, apply the approximation assertion on slightly
enlarged closed products still in the given domain and use
Cauchy estimates on the displayed products.
We join two adjacent \(g_i\)'s over
one transition, with control on its required face slices.

On a slightly enlarged \(D_b\times C\times P\), a chart spray
over \(f\), from Lemma~\ref{loc:spray}, expresses the two maps as
\[
 g_i=f^\sharp(\cdot,a_i),\qquad
 g_{i+1}=f^\sharp(\cdot,a_{i+1}),
\]
where \(a_i,a_{i+1}\) and their required derivatives are as small
as desired.  Choose a smooth function \(\chi:[0,1]\to[0,1]\)
equal to \(0\) near \(0\) and \(1\) near \(1\), and set
\[
 f_s=f^\sharp\bigl(\cdot,(1-\chi(s))a_i+\chi(s)a_{i+1}\bigr).
\]
This is a smooth \(s\)-family, holomorphic in \((t,y,u)\).
It equals the entire map \(g_i\) near \(s=0\) and \(g_{i+1}\)
near \(s=1\); on these subintervals use those entire definitions
on all \(y\).  It is arbitrarily close to \(f\) on the local
product, and its positive \(s\)-derivatives up to any fixed
order are arbitrarily small there.  The constants here depend
only on the previously fixed local chart and cutoffs.

The path $f_s$ is defined only over the small convex set $C$, except
near its endpoints, where it equals an entire map. We need a path
defined over the common large polydisc $B$, still close to $f_s$ on
$C^-$, and equal to the prescribed endpoint maps on all of $B$.
Direct interpolation in the chart over $f$ cannot provide this,
because that chart is not defined over $B$.

We obtain the larger domain by expanding the $y$ variable only
where the path already equals an entire endpoint map. After entire
approximation on the fixed small set $C$, we undo that expansion.
Choose \(d_*\in\operatorname{int}C\).  For some \(L\geq1\),
\begin{equation}\label{ana:eq:inverse-scaling}
 d_*+L^{-1}(B-d_*)\Subset\operatorname{int}C.
\end{equation}
Choose a smooth positive function \(\lambda:[0,1]\to[1,L]\)
which equals \(1\) outside the endpoint intervals where \(f_s\)
is an entire endpoint map, and equals \(L\) on smaller endpoint
intervals.  All of these choices are fixed.  The family
\[
 \widetilde f_s(t,y,u)
       =f_s\bigl(t,d_*+\lambda(s)(y-d_*),u\bigr),
 \qquad y\in C,
\]
is defined on a common neighborhood of
\([0,1]\times\overline D_b\times C\times P\).
This assertion uses entirety precisely where \(\lambda>1\);
where only the local definition of \(f_s\) is available,
\(\lambda=1\).
Apply Lemma~\ref{glue:interval} to make this family jointly
holomorphic near that compact product, with arbitrarily good
finite derivative accuracy.  Next apply
Lemma~\ref{ana:parameter-disc}, with \({\cal A}_{n-1}\) and
passive factors \(D_b,P\), to obtain an entire approximant
\(E(s,t,y,u)\).

For real \(s\), put
\[
 A_s(t,y,u)=
 E\bigl(s,t,d_*+\lambda(s)^{-1}(y-d_*),u\bigr).
\]
It is smooth in \(s\) and holomorphic in \((t,y,u)\) on the
whole prescribed large product.  On \(C'=C^-\Subset
\operatorname{int}C\), which contains all required overlap slices,
the inverse scaling stays in a fixed compact subset of
\(\operatorname{int}C\): by convexity it lies in the convex hull
of \(C'\cup\{d_*\}\). The two scalings cancel exactly:
\[
 \widetilde f_s\bigl(t,d_*+\lambda(s)^{-1}(y-d_*),u\bigr)
       =f_s(t,y,u).
\]
Thus approximating $\widetilde f_s$ on that fixed compact subset
makes \(A_s\) arbitrarily close to \(f_s\), with the required
derivatives, on these slices.
On the smaller endpoint intervals, \eqref{ana:eq:inverse-scaling}
gives the stronger statement that \(A_s\) is arbitrarily close
to \(g_i\), respectively \(g_{i+1}\), on the entire large
product \(\overline D_b\times B\times P\).
The chain rule costs only finite constants, since \(\lambda\)
and all scaling factors have already been fixed.

We may arrange exact endpoint joins.  A chart spray over
\(g_i\) exists on a slightly larger version of that large
product.  On a smaller initial \(s\)-interval, write
\(A_s=g_i^\sharp(\cdot,\beta_s)\).  Here \(\beta_s\), including
all required derivatives, is as small as desired.
Replace \(\beta_s\) by \(\kappa(s)\beta_s\), where \(\kappa\)
is \(0\) near \(0\) and \(1\) before leaving this endpoint
interval.  Do the analogous replacement near \(1\).
The repaired path, denoted \(H_s\), is exactly \(g_i\) near
one endpoint and exactly \(g_{i+1}\) near the other, and retains
the desired approximation on the overlap slices.  It need
only be holomorphic on a neighborhood of the fixed large
product, which the chart construction provides.

There is no circular smallness requirement in this procedure.
First choose the local \(g_i\)'s sufficiently accurate that
the interpolated \(f_s\) has the desired small error after the
fixed real transition cutoff.  Then, with \(g_i,\lambda,B\)
and the endpoint chart sprays fixed, choose the complexification
and entire approximation of \(\widetilde f_s\) as accurate as
needed to overcome all their finite derivative and chart
constants.  Finally perform the exact endpoint repairs.

Substitute a fixed smooth cutoff \(s=s(t)\) across the real
transition interval.  The result is a smooth real \(t\)-family,
holomorphic in \((y,u)\) on the large product, which joins the
two local entire maps and is as close to \(f\) as prescribed,
with all required tangential derivatives, on the face slices.
Doing this at the finitely many disjoint transition intervals
gives one such family on all of \(I\).  Extend it smoothly
past the endpoints of \(I\) using the first and last entire
maps \(g_i\), to which it already agrees on endpoint intervals.

Apply Lemma~\ref{glue:interval} once more, now to this real
\(t\)-family on the fixed large polydisc in \((y,u)\).
The result is jointly holomorphic near
\(I\times B\times P\) and retains the prescribed accuracy on
the face.  Proposition~\ref{ana:polydisc} is \({\cal A}_0\);
thus Lemma~\ref{ana:parameter-disc} with \(r=0\) converts this
last family into an entire map in \((t,y,u)\), still with the
required accuracy.
When \(n=1\), there are no \(y\) variables and the scaling
step is omitted; the same argument, or just this last
complexification and \({\cal A}_0\) step, applies.
\end{proof}

\subsection{Gluing across a real interface}

Lemma~\ref{ana:face} provides an entire upper-side map close to the
old map on the real face. Unlike the buffered overlap in
Lemma~\ref{glue:open}, the complex neighborhood on which this
accuracy persists may depend on the approximant. We therefore solve
the additive jump problem using only its H\"older trace. The
nonlinear correction will use the same quadratic iteration as
open-overlap gluing, now in this trace norm.

Fix $0<\alpha<1$. For a compact convex set $Q\subset\C^n$, write
\[
 Q_r=\{z:\operatorname{dist}(z,Q)<r\},\qquad
 M_r=Q_r\cap\{\operatorname{Im}z_1=0\},\qquad
 Q_r^\pm=Q_r\cap\{\pm\operatorname{Im}z_1>0\}.
\]
We always start with $r>0$, including when $Q$ has empty interior.
An auxiliary variable $v$ ranges over a ball $B_b\subset\C^N$.
All functions under consideration are holomorphic in $v$. All norms
below are uniform in $v$; the H\"older seminorm is taken in the
base variables only. Traces on $M_r$ are holomorphic in
$z'=(z_2,\ldots,z_n)$ on every available slice. For a side function set
\[
 \|a\|_{r,b,*}=\sup_{Q_r^\pm\times B_b}|a|
       +\sup_{v\in B_b}\|a(\,\cdot\,,v)|_{M_r}\|_{C^\alpha}.
\]
The H\"older norm on an open bounded set means the supremum norm plus
the supremum of all difference quotients in that set. The constants in
this section may depend on a bounded range of $r$, the dimensions,
$\alpha$, and the fixed joining geometry, but not on the input functions.

\begin{lemma}[Additive splitting at a real interface]\label{glue:linear-trace}
There are $C>0$ and an integer $p\geq1$ with the following property.
For $0<h<\min(r/2,1)$ and a vector-valued function $u$ on
$M_r\times B_b$ with the regularity just specified, there are linear
operators producing functions $a^\pm$, holomorphic on
$Q_{r-h}^\pm\times B_b$ and continuous up to $M_{r-h}\times B_b$, such that
\[
 a^+-a^-=u\quad\hbox{on }M_{r-h}\times B_b,
 \qquad
 \|a^+\|_{r-h,b,*}+\|a^-\|_{r-h,b,*}
       \leq Ch^{-p}\sup_{v\in B_b}\|u(\,\cdot\,,v)\|_{C^\alpha(M_r)}.
\]
The operators preserve holomorphic dependence on all auxiliary variables.
\end{lemma}

\begin{proof}
Use several intermediate thickenings between $Q_r$ and $Q_{r-h}$,
with successive distances fixed positive multiples of $h$. Cover the
required part of $M_r$ by coordinate boxes of side comparable to $h$.
The enlarged boxes used below are still inside $Q_r$. One can choose a
grid cover with bounded overlap, at most $C h^{-2n}$ boxes, and a smooth
subordinate partition whose derivatives of order $j$ are bounded by
$C_jh^{-j}$. Boxes outside a sufficiently large fixed ball are unnecessary.

In one enlarged box, with center $(t_0,z'_0)$ on the interface, take a
smooth real cutoff $\chi$ supported in its $t$-interval and equal to one
on a smaller interval. For $\operatorname{Im}\zeta\ne0$ put
\[
 C_u(\zeta,z',v)=\frac{1}{2\pi i}
       \int_\R\frac{\chi(t)u(t,z',v)}{t-\zeta}\,dt.
\]
The one-variable Cauchy boundary formula gives the jump $u$ on the
smaller interval, after assigning the signs of the two boundary values
in the usual way. Its H\"older bound follows by subtracting
$u(t_*,z',v)$ in the singular integral near a boundary point $t_*$:
the remaining numerator is bounded by a constant times
$|t-t_*|^\alpha$. The same subtraction in the difference of two boundary
values gives the $C^\alpha$ estimate. Rescaling the interval to unit
length gives a bound with a fixed power of $h^{-1}$.
The integral is holomorphic in $(z',v)$ by dominated integration.
Cauchy estimates on slightly smaller transverse boxes control transverse
derivatives, with a further fixed power of $h^{-1}$. Thus the local
side functions have the asserted joint H\"older trace bounds.

Denote these local pairs by $c_i^\pm$. Multiply by a smooth partition
$\theta_i$ in a band around the interface, equal in sum to one on a
smaller band, and extend the partition terms into the two sides with
support in their boxes. This gives smooth side functions
$f^\pm=\sum_i\theta_i c_i^\pm$ with trace jump $u$.
Their ordinary side derivatives $\bar\partial f^\pm$ fit together to a
single smooth form $\eta$ across the interface. Here is the point that
ensures smoothness despite the original H\"older data. In a smaller box
choose one reference pair $c_0^\pm$. Every difference
$c_i^\pm-c_0^\pm$ has equal traces, hence extends holomorphically across
the cut by continuity and Morera's theorem. Where $\sum_i\theta_i=1$,
\[
 \bar\partial f^\pm
    =\sum_i(\bar\partial\theta_i)(c_i^\pm-c_0^\pm).
\]
The right side is the restriction of a common smooth form.
Outside this smaller band, the terms are evaluated a distance comparable
to $h$ from their Cauchy singularities. Consequently any fixed number of
derivatives of $\eta$ is bounded by $C h^{-p_j}\|u\|_{C^\alpha}$.
It is closed, since it is locally a $\bar\partial$ derivative on the
open sides and smooth across their boundary.

Solve $\bar\partial w=\eta$ on an intermediate convex domain. A weighted
minimal $L^2$ solution with the fixed strictly plurisubharmonic weight
$|z|^2$ has $L^2$ norm at most $C\|\eta\|_{L^2}$ by the standard
$L^2$ theorem; the weight is bounded above and below on the domains in
question; this is the weighted estimate of
\cite[Theorem~2.2.1$'$]{HormanderLTwo}, also treated in
\cite{Hormander,Forstneric}. Interior estimates give both a
supremum and a first-derivative bound on the next smaller thickening,
with polynomial losses in $h^{-1}$. Indeed, writing
$\eta=\sum_{j=1}^n\eta_j\,d\bar z_j$, the distributional equation gives
$\Delta w=4\sum_j\partial\eta_j/\partial z_j$ componentwise.
Interior elliptic estimates on nested balls of radius comparable to
$h$, followed by Sobolev embedding, bound $w$ in $C^1$ by its $L^2$
norm and a fixed finite number of derivatives of $\eta$, with powers
of $h^{-1}$. This proves the required interior bound and hence the
H\"older trace bound for $w$.

The operator just used is fixed and linear on the base domain.
Applied to a holomorphic family of $L^2$ coefficients, it gives a
holomorphic family of $L^2$ solutions. Interior estimates turn this into
ordinary joint holomorphic dependence on $v$. Now set
$a^\pm=f^\pm-w$. Their side $\bar\partial$ derivatives vanish, their
trace jump is unchanged, and the stated estimate follows by enlarging
$p$ to dominate the finitely many derivative and covering losses.
\end{proof}

\begin{lemma}[Small nonlinear gluing on a real interface]\label{glue:trace}
The conclusion of Lemma~\ref{glue:open} holds for sprays on
$Q_r^-$ and $Q_r^+$ when their transition is
given on $M_r\times B_b$ and is arbitrarily small in
\[
 \|c\|_{r,b,\alpha}
       =\sup_{v\in B_b}\|c(\,\cdot\,,v)\|_{C^\alpha(M_r)}.
\]
Here the sprays are holomorphic on the open sides, continuous up to the
interface, and the transition is holomorphic in $(z',v)$ there. The
resulting glued map is holomorphic across the interface. No lower bound
is assumed for a complex normal thickness on which the input transition
is small.
\end{lemma}

\begin{proof}
Apply Lemma~\ref{glue:linear-trace} to split the transition error
as $b-a=c$. With losses $h$ in the base and fiber margins, its bound
is $Ch^{-p}e$, where $e=\|c\|_{r,b,\alpha}$; here and below the
side corrections are measured by their supremum and trace norms.
Cauchy estimates in the fiber variable give the same bound, with
additional fixed powers of $h^{-1}$, for their first and second
fiber derivatives on smaller balls.

As in Lemma~\ref{glue:open}, the corrected transition
$({\rm id}+b)^{-1}\circ\gamma\circ({\rm id}+a)={\rm id}+d$
is determined on the interface by
\[
 d=c(v+a)-c(v)-\{b(v+d)-b(v)\}.
\]
The $C^\alpha$ multiplication inequality and integration of the
first fiber derivative along line segments bound the first
difference by $Ch^{-p}e^2$. To compare its values at two base
points, the variation of the derivative at the shifted fiber
argument is controlled by the second fiber derivative. The same
estimates make the operation in braces contractive in the trace
norm. Thus, after increasing the fixed exponent $p$,
\[
                  \|d\|_{r-h,b-h,\alpha}\leq Ch^{-p}e^2.
\]
Only these finitely many derivatives are required at every step.
The losses $h_i=h_0 2^{-i}$ and the rescaled errors $E_i$ in the
proof of Lemma~\ref{glue:open} consequently give the same doubly
exponential convergence. The sums of side corrections and their
fiber derivatives converge on both sides, including their
$C^\alpha$ traces. Their limiting transitions are the identity.

Composing the original sprays with the limiting side maps gives
equal continuous traces. Near each interface point, continuity
places both images in a common target chart; Morera's theorem
across the real hyperplane then proves holomorphicity of the joined
map. The old-side approximation follows from the small fiber
corrections and the fixed spray's uniform continuity on compacta,
exactly as in Lemma~\ref{glue:open}.
\end{proof}

\subsection{Removal of convex faces}

The face lemma supplies an entire upper-side map with arbitrarily
accurate tangential derivatives along the join. It remains to attach
that map to the given lower-side map, retaining approximation on the
whole old convex set. Repeating this single-face step will reach a
polydisc, where Proposition~\ref{ana:polydisc} applies.

\begin{proof}[Proof of Theorem~\ref{ana:cap}]
We prove \({\cal A}_n\) by induction.  The case \(n=0\) is
Proposition~\ref{ana:polydisc}.  Suppose \({\cal A}_{n-1}\)
is known.

We first remove a single real affine constraint.  After an
invertible complex affine change of the \(n\) variables, write
the enlarged convex compactum as \(Q\subset\C^n\), and the
old compactum as
\[
 Q_- = Q\cap\{\operatorname{Im}z_1\leq0\}.
\]
Assume a map \(f\) is given near \(Q_-\times P\).
Cuts with no enlargement are trivial.  In the nontrivial
applications below, \(Q_-\) has nonempty interior and the
cut meets the interior of \(Q\). Figure~\ref{fig:real-interface}
shows the two sides to be joined.

\begin{figure}[ht]
\centering
\begin{tikzpicture}[x=0.9cm,y=0.9cm,>=Latex,font=\small]
 \begin{scope}
  \clip (0,0) ellipse (4.1 and 1.55);
  \fill[black!4] (-4.2,-1.6) rectangle (4.2,0);
  \fill[black!9] (-4.2,0) rectangle (4.2,1.6);
 \end{scope}
 \draw (0,0) ellipse (4.1 and 1.55);
 \draw[thick] (-4.1,0) -- (4.1,0);
 \node[fill=white,inner sep=3pt] at (0,0)
       {$F=Q\cap\{\operatorname{Im}z_1=0\}$};
 \node at (0,0.82) {$G$};
 \node at (0,-0.82) {$f^\sharp$ on $Q_-$};
 \node at (-3.6,1.38) {$Q$};
 \draw[->] (4.75,-1.35) -- (4.75,1.35);
 \node[above] at (4.75,1.35) {$\operatorname{Im}z_1$};
\end{tikzpicture}
\caption{A schematic real section of the single-face enlargement;
passive variables are suppressed. The lower region carries $f^\sharp$
and the upper region carries $G$. The input maps are compared only on
a slightly enlarged real face $F$, with no prescribed complex collar
and no bound for $G$ away from the interface.}
\label{fig:real-interface}
\end{figure}


Choose small convex thickenings of \(Q\) and of \(P\).
Their lower sides, together with additional small neighborhoods,
can be kept in the domain of \(f\).  To verify this point without
a transversality assumption on the boundary of \(Q\), note that
as \(\eta\downarrow0\),
\[
 (Q+\eta\overline B)\cap\{\operatorname{Im}z_1\leq0\}
       \longrightarrow Q_-
\]
in the following sufficient sense: every sequence of points
in the left-hand sets has all its limit points in \(Q_-\).
Compactness then places these sets in any prescribed
neighborhood of \(Q_-\), for small \(\eta\).
Here \(\overline B\) is the Euclidean closed unit ball.

Lemma~\ref{loc:spray} gives a chart spray
\(f^\sharp(z,u,v)\) over a neighborhood of the thickened
lower side, with a fixed auxiliary polydisc in \(v\).
Its trace on
\[
 \{\operatorname{Im}z_1=0\}
\]
is holomorphic near a convex compact face in
\((t,y)=(\operatorname{Re}z_1,z_2,\ldots,z_n)\).
Apply Lemma~\ref{ana:face}, treating \(u,v\) as passive
polydisc variables.  We obtain an entire map \(G(z,u,v)\)
arbitrarily close to \(f^\sharp\), with, say, two tangential
derivatives on a slightly enlarged face and with fixed
smaller \(u,v\) margins.

On this interface, invert the vertical chart of \(f^\sharp\).
For a fixed smaller \(v\)-polydisc this produces a transition
\[
 \gamma(z,u,v)=v+c(z,u,v)
\]
between \(G\) and \(f^\sharp\), with \(c\) as small as desired
in the \(C^\alpha\) trace norm, for any fixed \(0<\alpha<1\).
This follows from the just obtained tangential derivative
accuracy and the bounded derivatives of the fixed chart
inverse on compact sets; Cauchy estimates on the retained
\(u\)-margin also give the required H\"older control in \(u\).
The transition is holomorphic in the complex tangential variables
and in \(v\).
Use \(f^\sharp\) on the lower side and \(G\) on the upper
side. Apply Lemma~\ref{glue:trace} with convex base \(Q\times P\),
normal coordinate \(z_1\), and only \(v\) as the fiber variable.
With losses smaller than the chosen thickening margins, it glues
their suitably reparametrized
central maps on a neighborhood of \(Q\times P\).
Its corrections can be made arbitrarily small on the lower
side, so the resulting map approximates \(f\) on \(Q_-\times P\).
No estimate for \(G\) away from the interface is used.

Finally let \(K\subset\C^n\) be any nonempty compact convex set,
and let \(f\) be given near \(K\times P\).
Thicken \(K\) and \(P\) slightly within this domain, and choose
a full-dimensional real convex polytope \(K'\) containing \(K\)
in its interior and still within the available thickening.
Choose a large coordinate polydisc \(\overline\Delta_R^{\,n}\)
containing \(K'\) in its interior.  Write
\[
 K'=\bigcap_{j=1}^M\{\ell_j(z)\leq b_j\},
\]
where the \(\ell_j\)'s are real linear forms, and remove these
constraints one at a time inside
\(\overline\Delta_R^{\,n}\).
At each stage both compact sets are convex, the old one has
nonempty interior, and a nonredundant constraint can be put
in the form \(\operatorname{Im}z_1\leq0\) by a complex affine
coordinate change.  The single-constraint argument therefore
extends the map approximately to the next compactum.
There are only finitely many stages, so assign their errors
a sum smaller than half the desired tolerance.
The final map is defined near
\(\overline\Delta_R^{\,n}\times P\), a polydisc.
Proposition~\ref{ana:polydisc} approximates it by an entire map,
using the remaining half of the tolerance.
This proves \({\cal A}_n\), completes the induction, and gives
Theorem~\ref{ana:cap} by taking no passive variables.
\end{proof}
